\section{Obligation depth, open boxes, and exact depth two}
\label{sec:depth-two}

The preceding sections put the public trace of a sealed extension into normal form.  We now
ask the main semantic question: once that normal form has been reached, how far up the
coherence tower must the exact sealing obligation continue to grow?

There are two tempting but wrong answers.  The first says that unary action should be
enough.  This is false in a univalent setting, because an extension must respect exported
paths and equivalences, not merely exported objects.  The boolean swap obstruction below is
an explicit witness.  The second says that if binary comparison is needed, then an infinite
tower of higher primitive coherence fields should also be needed.  This is also false for the
fixed cubical sealed-extension calculus.  The higher structural obligations generated by
sealing are not arbitrary closed-boundary filler spaces.  They are open-box extension
problems whose missing face is part of the extension data.  Cubical composition and filling
supply a canonical extension.

Thus Section~\ref{sec:depth-two} proves a dichotomy:
\begin{quote}
Unary trace is too weak, binary trace is necessary, and every sealing-generated structural
obligation above binary trace is an open-box extension factor computed by the cubical core.
\end{quote}

The section is deliberately careful about levels.  Exact realized obligation objects may
still contain higher factors.  Those factors are contractible.  Minimal public signatures,
which count only irreducible primitive trace, do not contain those factors as primitive
fields.  Chronological locality is a third issue again: a binary field could in principle point
far back into history.  A separate factorization theorem below turns the arity-depth result
into the two-layer window used in the counting theorem.

\subsection{Exact obligations, primitive signatures, and historical arity}
\label{subsec:obligation-objects}

Let a candidate extension be sealed against a normalized historical interface.  For each
allowed depth bound we need both a syntactic signature of public obligations and the actual
semantic object whose inhabitants discharge those obligations.  The two should not be
confused.  Syntax determines what may be presented and counted; the realized object
determines whether the exact obligation space has stabilized.

\begin{definition}[Induced obligation signatures and realized objects]
\label{def:obsig-real}
For a candidate \(X\) sealed against a historical interface and for a depth parameter
\(k\), we write:
\[
  \mathrm{ObSig}_k(X),\qquad
  \mathrm{Ix}_k(X),\qquad
  \mathrm{Real}_k(X),\qquad
  \mathrm{Prim}_k(X),\qquad
  \mathrm{Der}_k(X).
\]
These have the following meanings.
\begin{center}
\begin{tabular}{p{0.18\linewidth}p{0.70\linewidth}}
\hline
\(\mathrm{ObSig}_k(X)\) & the finite normalized obligation signature visible up to
historical arity \(k\); \\
\(\mathrm{Ix}_k(X)\) & the finite set of primitive schema indices in
\(\mathrm{ObSig}_k(X)\); \\
\(\mathrm{Real}_k(X)\) & the dependent type, space, or cubical object interpreting
\(\mathrm{ObSig}_k(X)\); \\
\(\mathrm{Prim}_k(X)\) & the primitive irreducible sub-signature; \\
\(\mathrm{Der}_k(X)\) & the derived or transparently synthesized sub-signature. \\
\hline
\end{tabular}
\end{center}
The signature \(\mathrm{ObSig}_k(X)\) is read in the counting normal form of
Section~2: canonical telescope isomorphisms are quotiented, transparent reindexings are not
counted as new fields, and each field is tagged by its primitive or derived role.  The index
set \(\mathrm{Ix}_k(X)\) indexes only primitive fields.  The object
\(\mathrm{Real}_k(X)\), by contrast, is the exact semantic obligation object.
\end{definition}

The slogan is:
\[
  \mathrm{Real}_k(X) \text{ measures exact semantic obligation,}
  \qquad
  \mathrm{Prim}_k(X) \text{ measures primitive public syntax.}
\]
The main stabilization theorem is a theorem about \(\mathrm{Real}_k(X)\).  The trace-cost
bound is a theorem about \(\mathrm{Prim}_k(X)\) and about the trace part of normalized public
signatures.  These theorems interact, but neither one is merely a restatement of the other.

When a higher structural problem appears with one face still free, its exact semantic object
is represented by an extension type for the missing face together with a filler.  It is not
represented by the type of fillers for a completely fixed closed boundary.  This convention is
what makes the later open-box theorem applicable.

Because interfaces grow cumulatively, the primitive index sets and realized objects are
related by restriction maps
\[
  \mathrm{Ix}_1(X) \to \mathrm{Ix}_2(X) \to \mathrm{Ix}_3(X) \to \cdots,
  \qquad
  \mathrm{Real}_1(X) \to \mathrm{Real}_2(X) \to \mathrm{Real}_3(X) \to \cdots,
\]
where moving to the right exposes more historical interface.  The depth question asks when
these additional exposures stop changing the exact obligation object, up to canonical
equivalence.

\begin{remark}[Five statuses for a higher-looking field]
\label{rem:five-statuses}
The paper uses five different statuses that should not be collapsed.
\begin{center}
\small
\begin{tabular}{p{0.24\linewidth}p{0.66\linewidth}}
\hline
Status & Meaning \\
\hline
Present obligation type & A normalized obligation type genuinely occurs in some
\(\mathrm{Real}_k(X)\). \\
Contractible obligation type & The type occurs in \(\mathrm{Real}_k(X)\), but its fiber is
canonically contractible. \\
Transparently derived field & A field is part of a typing judgment or presentation, but is
uniformly reconstructed from exported lower data by cubical computation, transport, or
normalization. \\
Eliminated from a \(\mu\)-minimal signature & The field may appear in a presentation, but
not as a primitive field in the canonical normalized public signature measured by \(\mu\). \\
Absent from the exact obligation object & No normalized obligation type of that kind occurs
in the relevant \(\mathrm{Real}_k(X)\). \\
\hline
\end{tabular}
\end{center}
The higher horn factors considered below are present in exact obligation objects, are
contractible as open-box extension fibers, are transparently derived, and are eliminated from
\(\mu\)-minimal public signatures.  They are not absent from the exact semantics.
\end{remark}

\begin{definition}[Historical support and arity]
\label{def:historical-arity}
Let \(o\) be a normalized structural obligation.  Its support, written \(\mathrm{Supp}(o)\),
is the set of sealed layers whose atomic exported schemas occur irreducibly in the counting
normal form of the typing derivation of \(o\), after every prior layer has been collapsed to
its opaque normalized interface.  The historical arity of \(o\) is
\[
  a(o) := |\mathrm{Supp}(o)|.
\]
\end{definition}

Historical arity counts independent prior interfaces, not textual size.  A long formula that
mentions many parameters inside one opaque old signature can still be unary.  Conversely, a
short comparison can be binary if its typing derivation depends irreducibly on two
independent sealed exports.

\begin{lemma}[Arity-to-cell correspondence]
\label{lem:arity-cell}
In the structural-integration grammar of the fixed calculus, an irreducible obligation of
historical arity \(k\) has the shape of a \(k\)-dimensional structural coherence cell.  Arity
one gives action or transport data.  Arity two gives a naturality square or comparison
homotopy.  Arity \(k\geq 3\) gives a marginal higher horn-extension cell.
\end{lemma}

\begin{proof}
Unary action clauses mention one prior opaque signature, so their normalized obligations are
single action, substitution, or transport clauses.  Binary comparison clauses mention two
independent prior signatures and compare the two unary composites through them, giving a
square or homotopy.  The higher horn clauses are generated inductively: passing from arity
\(k-1\) to arity \(k\) exposes one further remote layer and asks for the missing comparison
face plus a filler of the resulting open box.  Since support is computed after opacity and
counting normalization, transparent aliases, record rebundlings, and transports along already
exported equalities cannot increase this arity.
\end{proof}

\begin{definition}[Obligation depth and minimal-signature depth]
\label{def:two-depths}
A fixed foundation or extension calculus has \emph{obligation depth} \(d_{\mathrm{obl}}\) if
\(d_{\mathrm{obl}}\) is the least integer such that, for every admissible candidate \(X\) and
all \(k\geq d_{\mathrm{obl}}\), there is a canonical equivalence
\[
  \mathrm{Real}_k(X) \simeq \mathrm{Real}_{d_{\mathrm{obl}}}(X).
\]
It has \emph{minimal-signature depth} \(d_\mu\) if \(d_\mu\) is the least integer such that
every admissible surface declaration \(e\) is presentation-equivalent to a declaration \(e'\)
whose primitive trace sub-signature contains no irreducible primitive historical field of
arity greater than \(d_\mu\).
\end{definition}

When this paper says simply ``coherence depth,'' it means obligation depth unless explicitly
qualified.  The minimal-signature depth is a public-accounting invariant.  It says which
fields remain primitive after normalization; it does not by itself say that exact higher
obligation objects are absent.

\begin{definition}[Chronological window]
\label{def:chronological-window}
A fixed extension calculus has chronological window size at most \(d\) if every admissible
sealing step at stage \(n+1\) has the following property after passage to the canonical
normalized public signature: every surviving primitive trace field has support contained in
\[
  \{L_n,L_{n-1},\ldots,L_{n-d+1}\},
\]
and every comparison with a remote layer \(L_{n-r}\), for \(r\geq d\), is
presentation-equivalent to a trace term factored through the exported normalized signatures of
those \(d\) recent layers.
\end{definition}

A bound on arity is not yet a chronological bound.  An arity-two comparison could compare a
new layer directly with a very old one.  The recent-history theorem below supplies the extra
factorization needed before the recurrence theorem in Section~6 can count only a two-layer
window.

\begin{remark}[Level audit]
\label{rem:level-audit}
The main results of this section live at different levels.
\begin{center}
\small
\begin{tabular}{p{0.29\linewidth}p{0.29\linewidth}p{0.31\linewidth}}
\hline
Result type & What it bounds & What it does not yet prove \\
\hline
\(\mathrm{Real}_k(X)\simeq \mathrm{Real}_2(X)\) & exact realized obligation objects &
which surface fields are primitive in a public presentation \\
\(d_\mu\leq 2\) & primitive field arity in \(\mu\)-minimal signatures & a two-layer
chronological window \\
Recent-history factorization & primitive chronological support & the arithmetic recurrence
without a coupling hypothesis \\
Full-coupling recurrence & primitive trace cardinality & the cubical depth theorem itself \\
\hline
\end{tabular}
\end{center}
The proof below uses the first three rows.  The fourth row is the subject of the next
section.
\end{remark}

\subsection{A contrast case: UIP collapses exact obligation depth to one}
\label{subsec:UIP-depth-one}

Before proving the cubical theorem, it is useful to see what happens in a setting where
higher identity structure is proof-irrelevant.  In extensional or UIP-like systems the
coherence problem really does collapse much further.

\paragraph{What this contrast says.}
In a UIP setting, higher exact equality obligations carry no independent information once
their endpoints are fixed.  Obligation depth collapses to one.

\paragraph{What it does not say.}
It is not the cubical theorem.  It does not describe univalent binary trace, and it does not
justify ignoring binary naturality in cubical foundations.

\paragraph{Where the assumption enters.}
The collapse uses UIP directly: any two identity proofs between fixed endpoints are equal.

\begin{theorem}[UIP collapse]
\label{thm:UIP-depth-one}
In Martin-Lof type theory with UIP, the obligation depth is \(1\).
\end{theorem}

\begin{proof}
Under UIP, every identity type is a set in the relevant sense, and every higher identity type
between fixed endpoints is a proposition.  Once such a higher identity type is inhabited, it
is contractible.  By Lemma~\ref{lem:arity-cell}, every normalized coherence obligation of
dimension at least two corresponds to a higher identity witness over a boundary already fixed
by lower data.  UIP makes that witness unique.  Thus every obligation of arity at least two
is determined by its unary boundary, and
\[
  \mathrm{Real}_k(X) \simeq \mathrm{Real}_1(X)
  \qquad (k\geq 1).
\]
Depth zero would mean that a nontrivial sealed extension need not interact with the current
interface at all, so the least stabilizing depth is one.
\end{proof}

This contrast helps locate the cubical result.  Cubical univalence prevents the same collapse
at depth one; exported equivalences create genuine binary comparison obligations.  But Kan
composition and filling still prevent the sealing-generated structural tower from producing
primitive public trace above depth two.

\subsection{The cubical upper bound: higher structural obligations are open boxes}
\label{subsec:cubical-upper-bound}

We now specialize to the fixed sealed-extension calculus over the chosen cubical core.  The
ambient core has interval structure, transport, partial elements, homogeneous and dependent
composition, filling, and univalence.  Full coupling is not used in this subsection.  The
upper bound depends on the structural horn grammar and on cubical open-box computation;
full coupling returns only in Section~6, when the two-layer window is counted.

The important semantic clarification is this: the paper does not claim that arbitrary higher
filler spaces in a univalent type theory are contractible.  They are not.  The relevant object
is more specific.  It is the theorem-facing extension type of a structural open box whose
visible faces are fixed by prior public trace and whose missing remote face is still part of
the extension data.

\begin{definition}[Structural horn-extension package]
\label{def:structural-horn-package}
Let a candidate \(X\) be sealed at stage \(n+1\), and let \(k\geq 3\).  The remote layer
exposed by the depth-\(k\) step is
\[
  R_k(n) := L_{n-k+1}.
\]
For a boundary \(\partial : \mathrm{Real}_{k-1}(X)\), the visible boundary object
\[
  \partial_{\mathrm{vis}}^{(k)}(X,R_k(n))
\]
consists of the already exported unary traces of \(X\) against the recent layers, the binary
comparison traces exported by the intervening layers, and all degenerate, reindexed, or
transport faces generated transparently from those public traces.

The only face not included in this visible boundary is the direct remote comparison between
the action of \(X\) and the newly exposed layer \(R_k(n)\).  The telescope of such missing
comparison faces compatible with the endpoints of \(\partial_{\mathrm{vis}}^{(k)}\) is written
\[
  \Gamma_{\mathrm{miss}}^{(k)}(\partial).
\]
For \(\gamma:\Gamma_{\mathrm{miss}}^{(k)}(\partial)\), the filler predicate
\[
  \mathrm{Fill}^{(k)}_{\partial}(\gamma)
\]
is the type of cubical fillers completing the open \(k\)-box whose visible faces are
\(\partial_{\mathrm{vis}}^{(k)}\) and whose missing face is \(\gamma\).  The structural
horn-extension object is
\[
  H_k(\partial)
  :=
  \sum_{\gamma:\Gamma_{\mathrm{miss}}^{(k)}(\partial)}
  \mathrm{Fill}^{(k)}_{\partial}(\gamma).
\]
\end{definition}

This is the central formal move of the paper.  The marginal higher structural datum is a
missing-face-plus-filler package.  It is not a closed-boundary filler type.  The missing face is
chosen at the same time as the filler, and cubical composition computes the canonical choice.

\begin{definition}[Open-box extension type]
\label{def:open-ext}
Let \(\Gamma\) be a cubical context, let \(A : \Gamma,t:\mathbb I\vdash \mathcal U\) be a
fibrant family, let \(\varphi\) be a visible-face cofibration, and suppose given compatible
open-box data
\[
  u : (t:\mathbb I) \to \mathrm{Partial}_{\varphi}(A(t))
\]
and
\[
  u_0 : A(i0)[\varphi|_{t=i0}\mapsto u(i0)].
\]
The open-box extension type is
\[
  \mathrm{OpenExt}_{\Gamma}(A,\varphi,u,u_0)
  :=
  \sum_{\ell:A(i1)[\varphi|_{t=i1}\mapsto u(i1)]}
  \mathrm{Fill}(A,\varphi,u,u_0,\ell),
\]
where \(\mathrm{Fill}(A,\varphi,u,u_0,\ell)\) is the type of relative sections
\(f:(t:\mathbb I)\to A(t)\) satisfying the base equation at \(i0\), the lid equation at
\(i1\), and agreement with the partial boundary \(u(t)\) on \(\varphi\).
\end{definition}

The lid \(\ell\) is the missing remote face.  The second component is the filler.  Taking
both components together is what turns the open-box extension problem into a contractible
fiber.

\begin{lemma}[Missing-face packages are open-box extensions]
\label{lem:horn-open-box}
For every structural horn generated by the sealing grammar, the missing-face package
\(H_k(\partial)\) is canonically equivalent to an open-box extension type
\[
  \mathrm{OpenExt}_{\Gamma}(A_k,\varphi_k,u_k,u^0_k).
\]
The equivalence is natural in cubical substitution.
\end{lemma}

\begin{proof}
Expand \(\Gamma_{\mathrm{miss}}^{(k)}(\partial)\) as the omitted remote face of the structural
box.  A choice of \(\gamma\) together with a filler is exactly a lid at the top face agreeing
with the visible partial boundary, together with a relative filler from the base face to that
lid.  Reassociating this dependent sum gives the open-box extension type of
Definition~\ref{def:open-ext}.  Restriction to faces, reassociation of dependent sums, and
reindexing of partial elements all commute with cubical substitution, giving naturality.
\end{proof}

\begin{theorem}[Structural open-box extension contractibility]
\label{thm:open-ext-contr}
For every structural open box \((A,\varphi,u,u_0)\) generated by the sealing grammar in the
fixed cubical core,
\[
  \Gamma \vdash
  \mathrm{isContr}\bigl(\mathrm{OpenExt}_{\Gamma}(A,\varphi,u,u_0)\bigr).
\]
A center is computed by cubical composition and filling:
\[
  \bigl(\mathrm{comp}(A,\varphi,u,u_0),
        \mathrm{hfill}(A,\varphi,u,u_0)\bigr).
\]
Moreover, for every cubical substitution \(\sigma:\Delta\to\Gamma\), there is a canonical
equivalence
\[
  \mathrm{OpenExt}_{\Gamma}(A,\varphi,u,u_0)[\sigma]
  \simeq
  \mathrm{OpenExt}_{\Delta}(A[\sigma],\varphi[\sigma],u[\sigma],u_0[\sigma])
\]
that maps centers to centers, with contraction paths stable under substitution.
\end{theorem}

\begin{proof}
The visible boundary is encoded as a cubical partial element.  The compatibility equations on
overlaps are exactly the well-typedness conditions for \(u\) and for the subtype element
\(u_0\).  Cubical composition computes the top lid; cubical filling supplies the relative
section from the base to that lid.  This gives the displayed center.

Let \((\ell,f)\) be any other lid-and-filler pair.  To contract it to the canonical pair, apply
cubical composition one dimension higher in the extension family.  The boundary of this
higher composition is given by the candidate filler \(f\), the canonical filler, the visible
partial boundary on \(\varphi\), and the common base face at \(t=i0\).  The endpoint and
overlap equations required for this higher box are precisely the filler laws for \((\ell,f)\)
and for the canonical extension.  The resulting higher filler is a path from \((\ell,f)\) to
the canonical extension.

Uniformity of composition and filling under substitution gives the displayed reindexing
equivalence and center-stability statement.  No presentation quotient is used in this proof;
it is a theorem about the theorem-facing cubical extension object.
\end{proof}

\begin{lemma}[Substitution stability of structural open boxes]
\label{lem:open-box-substitution}
For every substitution \(\sigma:\Delta\to\Gamma\), the reindexed tuple
\((A_k,\varphi_k,u_k,u^0_k)[\sigma]\) is the structural open box generated by the same
sealing rule in context \(\Delta\).  The canonical center and contraction of
Theorem~\ref{thm:open-ext-contr} are stable under this reindexing.
\end{lemma}

\begin{proof}
Substitution acts structurally on interval expressions, face formulas, partial elements,
subtypes, and the lower public boundary data from which \((A_k,\varphi_k,u_k,u^0_k)\) is
assembled.  The cubical composition and filling operations are uniform under substitution.
Therefore the reindexed open box is exactly the open box generated after first substituting
the sealing rule, and the computed lid, filler, and contraction reindex to the corresponding
computed lid, filler, and contraction in \(\Delta\).
\end{proof}

\begin{remark}[Open boxes are not closed boundaries]
\label{rem:open-not-closed}
Theorem~\ref{thm:open-ext-contr} is an open-box extension theorem.  It does not assert that
the type of fillers for an arbitrary closed boundary is contractible.  For a closed boundary,
the filler fiber can carry genuine homotopy.  The contractibility used here comes from the
fact that the missing remote face is part of the extension data.  In slogan form:
\[
  \text{open-box extension type: contractible;}
  \qquad
  \text{arbitrary closed-boundary filler type: not claimed.}
\]
This distinction is the reason the paper models marginal higher structural obligations as
missing-face-plus-filler packages.
\end{remark}

\begin{lemma}[Syntactic horn reduction for structural integration]
\label{lem:syntactic-horn-reduction}
Let \(X\) be a foundational-core candidate sealed at stage \(n+1\) in the fixed cubical
extension calculus.  For every \(k\geq 3\) and every boundary
\(\partial:\mathrm{Real}_{k-1}(X)\), the additional data exposed by adjoining the remote layer
\(L_{n-k+1}\) are canonically represented by the horn-extension fiber
\[
  H_k(\partial)
  :=
  \sum_{\gamma:\Gamma_{\mathrm{miss}}^{(k)}(\partial)}
  \mathrm{Fill}^{(k)}_{\partial}(\gamma).
\]
Every face of the box except the direct remote comparison face is already exported by unary
and binary trace of the intervening sealed layers, together with transparent transport and
degenerate faces.
\end{lemma}

\begin{proof}
By construction, \(\mathrm{Real}_k(X)\) records only structural integration obligations induced
by sealing \(X\) against opaque prior layers.  Earlier layers therefore contribute only their
exported payload and trace schemas; arbitrary higher-algebraic structure internal to those
layers is invisible at this sealing boundary.  Passing from depth \(k-1\) to depth \(k\)
exposes exactly one further remote layer, namely \(L_{n-k+1}\).

The intervening sealed layers have already exported the unary bridges and binary comparison
traces linking that remote layer to more recent history.  Counting normalization treats the
associated reindexings, degenerate faces, and transports as transparent.  Thus all visible
faces of the comparison box are fixed by lower public trace.  The only new component is the
direct remote comparison face \(\gamma\), and once \(\gamma\) is chosen the remaining exact
obligation is precisely the filler predicate.  This is the displayed dependent sum.
\end{proof}

\paragraph{A concrete depth-three picture.}
At depth three, the new layer \(X\) is compared with \(L_n\), \(L_{n-1}\), and the newly
exposed remote layer \(L_{n-2}\).  The visible faces include the unary bridges from \(X\) to
the recent layers and the binary comparisons already exported by \(L_n\) and \(L_{n-1}\).
The missing face is the direct comparison between the two composites from \(X\) to
\(L_{n-2}\).  The depth-three factor is therefore
\[
  H_3(\partial)
  =
  \sum_{\gamma:\Gamma_{\mathrm{miss}}^{(3)}(\partial)}
  \mathrm{Fill}^{(3)}_{\partial}(\gamma).
\]
The term \(\gamma\) is not new payload.  It is the missing remote face of the structural open
box.  Cubical composition computes a canonical \(\gamma\), and filling packages the
accompanying three-dimensional cell.

\subsection{Derivedness and minimal public signatures}
\label{subsec:derivedness-minimal-signatures}

The open-box theorem is semantic.  It says that the higher exact factor is contractible.  The
public-signature theorem is syntactic-accounting.  It says that an explicit higher structural
field in a surface presentation can be replaced by the cubically computed term and therefore
removed from the primitive part of a \(\mu\)-minimal signature.

The bridge between the two levels is the computational replacement theorem from
Section~3: when a presented field is accompanied by a uniform derivation from lower public
boundary data, the presentation relation allows the field to be replaced by that derived term.
The open-box theorem supplies exactly such a derivation for higher structural trace.

\begin{theorem}[Minimal-signature elimination above binary arity]
\label{thm:minimal-signature-elimination}
Let \(e\) be an admissible surface extension declaration over \(B_n\).  Then there exists a
presentation-equivalent declaration \(e'\approx e\) whose normalized trace signature contains
no irreducible primitive historical field of arity \(k\geq 3\).  Equivalently, every canonical
\(\mu\)-minimal normalized public signature contains only unary and binary primitive
structural trace fields.

More concretely, if a presented structural field of arity \(k\geq 3\) occurs in some surface
declaration, then in the canonical trace-presentation interface it is transparently derived
from its lower-arity boundary by a term synthesized from cubical composition, transport, and
filling.
\end{theorem}

\begin{proof}
Fix a presentation of \(e\) and an apparent primitive structural trace field \(\phi\) of arity
\(k\geq 3\) in its elaborated sealed export.  By the normal-form classification of Section~4,
\(\phi\) is a higher structural horn package rather than payload.  Let \(\partial\) be its
lower-arity boundary data.

By Lemma~\ref{lem:syntactic-horn-reduction}, the additional datum contributed by this
arity-\(k\) step is a horn-extension fiber \(H_k(\partial)\).  By
Lemma~\ref{lem:horn-open-box}, this fiber is the corresponding open-box extension type.
By Theorem~\ref{thm:open-ext-contr}, the open-box extension has a canonical center computed
from \(\partial\) by the cubical operations.  Write this computed term as
\[
  t_{\phi}(\partial)
  :=
  \bigl(\mathrm{comp}(\partial),\mathrm{hfill}(\partial)\bigr).
\]
This term has the same boundary behavior and the same structural role as the presented
field \(\phi\).

The computational replacement rule of Section~3 then applies: a presented higher structural
field equipped with such a uniform lower-boundary derivation is presentation-equivalent to
the declaration in which the field is replaced by \(t_{\phi}(\partial)\) and marked derived.
Repeating this replacement for every arity-\(k\geq 3\) structural field terminates because the
normalized signature is finite.  The resulting presentation-equivalent declaration \(e'\) has
no primitive structural field above binary arity.  Since \(\mu\) counts only primitive trace
fields in canonical normal form, every \(\mu\)-minimal signature has the same property.
\end{proof}

\begin{remark}[What is eliminated]
Theorem~\ref{thm:minimal-signature-elimination} eliminates primitive public fields, not exact
semantic factors.  A higher horn factor may still occur in \(\mathrm{Real}_k(X)\), but after
Theorem~\ref{thm:open-ext-contr} it occurs as a contractible factor with a canonical center.
The trace-cost invariant \(\mu\) does not count such a factor as primitive public trace.
\end{remark}

\subsection{Exact stabilization above binary trace}
\label{subsec:exact-stabilization}

We can now prove the upper bound on obligation depth.  The proof is conceptually simple:
from depth two onward, each further exposure of remote history adjoins one structural
open-box extension factor; each such factor is contractible; dependent sums with
contractible fibers are equivalent to their bases.

\begin{lemma}[Kan contraction for structural horn fibers]
\label{lem:kan-contraction-horn-fibers}
For every structural horn-extension fiber \(H_k(\partial)\) generated by the fixed cubical
sealing grammar, the projection
\[
  H_k(\partial) \to 1
\]
is contractible in the cubical metatheory.  The center is the missing face and filler computed
from \(\partial\) by the cubical operations.
\end{lemma}

\begin{proof}
By Lemma~\ref{lem:horn-open-box}, \(H_k(\partial)\) is the open-box extension type for the
structural tuple \((A_k,\varphi_k,u_k,u^0_k)\).  Applying
Theorem~\ref{thm:open-ext-contr} gives a canonical center and a contraction of the whole
extension space.  Presentation equivalence is not used to prove this contractibility; it is
used only in Theorem~\ref{thm:minimal-signature-elimination} to classify the corresponding
public field as derived.
\end{proof}

\begin{theorem}[Exact stabilization of realized structural obligations]
\label{thm:exact-stabilization}
In the fixed cubical extension calculus, for every foundational-core candidate \(X\) and
every \(k\geq 2\) there is a canonical equivalence
\[
  \mathrm{Real}_k(X) \simeq \mathrm{Real}_2(X).
\]
Consequently \(d_{\mathrm{obl}}\leq 2\).
\end{theorem}

\begin{proof}
The case \(k=2\) is immediate.  For \(k=3\), Lemmas~\ref{lem:arity-cell} and
\ref{lem:syntactic-horn-reduction} identify the marginal depth-three datum over each
boundary \(\partial:\mathrm{Real}_2(X)\) as the horn-extension fiber \(H_3(\partial)\).  Hence
\[
  \mathrm{Real}_3(X)
  \simeq
  \sum_{\partial:\mathrm{Real}_2(X)} H_3(\partial).
\]
By Lemma~\ref{lem:kan-contraction-horn-fibers}, every fiber \(H_3(\partial)\) is
contractible.  A dependent sum over a family of contractible fibers is canonically equivalent
to its base, so
\[
  \mathrm{Real}_3(X) \simeq \mathrm{Real}_2(X).
\]

For \(k>3\), the same argument repeats.  Passing from depth \(k-1\) to depth \(k\) exposes
one additional remote layer and adjoins one structural horn-extension fiber \(H_k(\partial)\)
over \(\partial:\mathrm{Real}_{k-1}(X)\).  Lemma~\ref{lem:kan-contraction-horn-fibers}
contracts that fiber, so
\[
  \mathrm{Real}_k(X) \simeq \mathrm{Real}_{k-1}(X)
  \qquad (k\geq 3).
\]
Iterating these equivalences yields \(\mathrm{Real}_k(X)\simeq\mathrm{Real}_2(X)\) for all
\(k\geq 2\).
\end{proof}

\begin{corollary}[Contractible-factor decomposition]
\label{cor:contractible-factor-decomposition}
For every foundational-core candidate \(X\) and every \(k\geq 2\), there is a canonical
equivalence
\[
  \mathrm{Real}_k(X)
  \simeq
  \mathrm{Real}_2(X) \times C_k(X),
\]
where \(C_2(X):=1\) and, for \(k>2\), \(C_k(X)\) is the iterated chain of canonical remote
horn factors selected by the cubical compositions in the proof of
Theorem~\ref{thm:exact-stabilization}.  Each \(C_k(X)\) is canonically contractible.
\end{corollary}

\begin{proof}
The proof of Theorem~\ref{thm:exact-stabilization} trivializes each successive horn-extension
fiber by choosing its canonical cubical center.  Iterating those choices identifies the total
higher part with a product of the depth-two boundary and the chain of chosen remote horn
factors.  Each factor is contractible by Lemma~\ref{lem:kan-contraction-horn-fibers}, hence
the whole chain \(C_k(X)\) is contractible.
\end{proof}

The corollary expresses the exact semantic situation more vividly than the bare equivalence:
above depth two there may be additional exact factors, but all of them are canonical and
contractible.  The primitive public signature sees only the depth-two part.

\subsection{From arity depth to a two-layer chronological window}
\label{subsec:recent-history}

The exact-depth theorem bounds arity.  It does not, by itself, identify which layers must be
inspected in time.  A primitive binary field could still compare a new layer with a remote
old layer.  The next theorem rules out such primitive remote binary fields for the fixed
calculus, using the factorization-complete public exports introduced earlier.

\begin{theorem}[Recent-history factorization]
\label{thm:recent-history-factorization}
In the fixed cubical extension calculus, every binary trace field that survives in a
\(\mu\)-minimal normalized public signature at stage \(n+1\) factors through the previous two
normalized signatures.  More precisely, let \(X\) be sealed at stage \(n+1\).  Once the unary
and binary trace data of \(X\) against \(L_n\) and \(L_{n-1}\) have been fixed, every
comparison between \(X\) and a remote layer \(L_{n-r}\), for \(r\geq 2\), is
definitionally equal after normalization to a term obtained by iterating the already exported
unary bridges and binary comparison traces of the intervening normalized signatures.
Consequently no \(\mu\)-minimal public signature contains a primitive remote binary field.
\end{theorem}

\begin{proof}
Consider first the layer \(L_{n-2}\).  The sealed export of \(L_n\) contains the adjacent
unary bridge to \(L_{n-1}\), and the sealed export of \(L_{n-1}\) contains the adjacent bridge
to \(L_{n-2}\).  The binary part of the trace exported by these layers records the comparison
between the adjacent composite and the corresponding direct public route.  Since exports are
factorization-complete, the degenerate faces and transport faces needed to assemble the next
comparison box are transparent public data.

When \(X\) is sealed against \(L_n\) and \(L_{n-1}\), it supplies the unary comparisons to
those two layers and the binary square comparing the two routes through them.  Combining
this data with the exported adjacent bridges determines all visible faces of an open
three-box.  The missing face is exactly the would-be direct comparison from \(X\) to
\(L_{n-2}\).  Thus the remote comparison is not an independent primitive binary field; it is
the lid of a structural horn-extension object.  Cubical composition computes that lid, and
cubical filling supplies the accompanying three-cell.  By Theorem~\ref{thm:minimal-signature-elimination},
the presented remote comparison is therefore replaced by the derived term in every
\(\mu\)-minimal signature.

The argument iterates.  Suppose the comparison from \(X\) to \(L_{n-r+1}\) has already
been factored through recent exported signatures.  The normalized export of \(L_{n-r+1}\)
contains the adjacent bridge to \(L_{n-r}\), the binary comparison trace for the composite
through that bridge, and the transparent degenerate faces needed for the next open box.
Combining these with the already factored comparison for \(X\) produces another structural
horn-extension problem whose missing face is the comparison from \(X\) to \(L_{n-r}\).
Cubical composition and filling synthesize that face and filler.  Hence every remote binary
comparison is definitionally generated from the recent public exports and is not primitive in
canonical minimal form.
\end{proof}

\begin{corollary}[Two-layer chronological window]
\label{cor:two-layer-window}
In the fixed cubical extension calculus, every sealing step at stage \(n+1\) has chronological
window size at most two.  All primitive obligations are generated by the exported traces of
\(L_n\) and \(L_{n-1}\), while comparisons with layers \(L_{n-r}\) for \(r\geq 2\) factor
through those two normalized signatures and are synthesized by iterated cubical filling.
\end{corollary}

\begin{proof}
Theorem~\ref{thm:exact-stabilization} shows exact arity stabilization at depth two, while
Theorem~\ref{thm:recent-history-factorization} shows that surviving binary primitive fields
can be chosen to factor through the two most recent normalized signatures.  Therefore deeper
history contributes only derived fillers and already mediated boundary data, not primitive
public trace fields outside the recent two-layer window.
\end{proof}

The recurrence theorem in Section~6 starts only after this point.  It counts the active basis
inside the two-layer chronological window under an additional full-coupling hypothesis.

\subsection{The lower bound: univalence forces binary trace}
\label{subsec:swap-lower-bound}

It remains to show that the upper bound is exact.  The obstruction is small but decisive.
A univalent interface that exports the two-point type also exports the path corresponding to
its nontrivial automorphism.  A sealed action must respect that path.  Objectwise unary data
alone cannot express this naturality condition.

\begin{theorem}[Explicit binary sealing obstruction]
\label{thm:swap-obstruction}
Let the active interface of a univalent foundation export the two-point type \(2\) and the
univalent path
\[
  p_{\mathrm{swap}} := \mathrm{ua}(\mathrm{swap}) : 2 = 2.
\]
Then there exists an objectwise promoted-interface candidate whose unary clause is
well-formed on the active object \(2\) but whose sealing is not determined by unary data.
Sealing induces a genuine binary obligation over \(p_{\mathrm{swap}}\).  Consequently depth-one
data are insufficient in general, and \(d_{\mathrm{obl}}\geq 2\).
\end{theorem}

\begin{proof}
Consider an active-interface package consisting of an endomap on the exported two-point
type:
\[
  X_{\mathrm{act}} := (F_2 : 2\to 2).
\]
The unary candidate chooses the constant-left map
\[
  F_2 := \lambda\_.\,\mathsf{left}.
\]
This is a well-formed objectwise clause.  But sealing a globally acting structural extension
requires compatibility with exported paths.  Since the interface exports
\(p_{\mathrm{swap}}\), the binary naturality obligation is
\[
  \alpha_{\mathrm{swap}}:
  \mathrm{transport}_{\lambda A. A\to A}(p_{\mathrm{swap}},F_2)=F_2.
\]
By univalence, transport in the endomorphism family along \(p_{\mathrm{swap}}\) computes by
conjugation with the swap equivalence.  Therefore
\[
  \mathrm{transport}_{\lambda A. A\to A}(p_{\mathrm{swap}},F_2)
  =
  \mathrm{swap}\circ F_2\circ \mathrm{swap}^{-1}
  =
  \lambda\_.\,\mathsf{right}.
\]
Thus \(\alpha_{\mathrm{swap}}\) would be an equality
\[
  (\lambda\_.\,\mathsf{right})=(\lambda\_.\,\mathsf{left}),
\]
which is impossible because evaluating at \(\mathsf{left}\) would force
\(\mathsf{right}=\mathsf{left}\).

There is also a positive control: if \(F_2\) is the identity map, conjugation along swap
returns the identity map, and the binary obligation is inhabited.  The failure of the
constant-left map is therefore not a defect of the formulation; it is the type checker
detecting an objectwise action that does not respect exported path structure.  Hence
\(\mathrm{Real}_2(X)\) can contain a binary obligation not present in \(\mathrm{Real}_1(X)\),
and the obligation depth cannot be one.
\end{proof}

\begin{remark}[Adjunctions as an explanatory analogue]
Adjunctions exhibit the same binary pattern categorically.  Once the object action and the
unit/counit data have been fixed, the triangle identities compare composites already
determined at the unary level.  The paper uses this only as an analogy.  The internal lower
bound for the fixed cubical calculus is the swap obstruction of
Theorem~\ref{thm:swap-obstruction}.
\end{remark}

\subsection{A topological example: clutching as unary/binary separation}
\label{subsec:clutching-example}

The swap obstruction is the load-bearing lower bound.  It is nevertheless useful to record a
more geometric example of the same separation between object data and comparison data.  The
clutching construction for bundles over the two-sphere supplies such a picture.

\begin{remark}[Clutching family example]
\label{rem:clutching}
The higher-inductive presentation of the two-sphere as a suspension has point constructors
\[
  N,S:S^2
\]
and meridian paths
\[
  \mathsf{merid}:S^1\to (N=S).
\]
To specify a type family \(P:S^2\to\mathcal U\), the dependent eliminator asks for pole
fibers
\[
  P_N:\mathcal U,
  \qquad
  P_S:\mathcal U,
\]
and a path family
\[
  \gamma: \prod_{x:S^1} P_N=P_S.
\]
By univalence, this path family can be read as a family of equivalences between the two pole
fibers.  Since \(S^1\) has a point and a loop, specifying \(\gamma\) contains a unary-looking
base equivalence together with a binary comparison over the loop.  The loop comparison is
genuine comparison data; it cannot be recovered from the endpoints alone.

This example illustrates the same depth-two pattern as the swap obstruction: binary
comparison data are real.  But it is not used to prove the universal upper bound.  The upper
bound comes from the structural horn-to-open-box theorem above.  In this particular HIT
presentation, once the binary clutching data are supplied, further compatibility conditions
appear as horn-extension problems over the displayed boundary and are therefore derived in
the fixed cubical trace calculus.
\end{remark}

\subsection{Exact public structural depth two}
\label{subsec:exact-depth-two-conclusion}

We can now combine the upper and lower bounds.

\begin{corollary}[Fixed-calculus exact structural integration depth two]
\label{cor:exact-depth-two}
In the fixed cubical sealed-extension calculus, admissible sealed structural extensions have
exact structural integration depth two:
\[
  d_{\mathrm{obl}}=2.
\]
Moreover every \(\mu\)-minimal normalized public signature contains no primitive structural
trace field of historical arity greater than two, and every sealing step has a two-layer
chronological window under the factorization-complete export hypothesis.
\end{corollary}

\begin{proof}
Theorem~\ref{thm:exact-stabilization} gives the upper bound
\(d_{\mathrm{obl}}\leq 2\) by showing
\(\mathrm{Real}_k(X)\simeq\mathrm{Real}_2(X)\) for all \(k\geq 2\).  Theorem~\ref{thm:minimal-signature-elimination}
turns the same cubical horn computation into the public-signature statement
\(d_\mu\leq 2\).  Theorem~\ref{thm:recent-history-factorization} and
Corollary~\ref{cor:two-layer-window} give the separate chronological factorization needed
for the later recurrence.  Finally, Theorem~\ref{thm:swap-obstruction} gives the lower bound
\(d_{\mathrm{obl}}\geq 2\) by exhibiting a genuine binary sealing obstruction in the univalent
setting.  Hence the exact obligation depth is two.
\end{proof}

The result should be read with the distinctions of Remark~\ref{rem:five-statuses} in mind.
Higher structural exact factors do not vanish from the semantics; they stabilize as
contractible open-box extension factors.  What disappears from the primitive public
interface is the need to present those factors as irreducible fields.  This is the semantic
content behind the paper's slogan:
\[
  \text{binary trace is necessary; higher structural trace is cubically payable.}
\]
